\begin{frame}[fragile]{Installation of UHD from Source Code (1/5)}

\footnotesize

\begin{itemize}
    \item \textbf{UHD can be installed in two ways:}
    \begin{itemize}
        \item \textbf{Binary Installation:} Installs pre-built UHD packages using the operating system package manager (e.g., \texttt{apt}). It is simple and quick, but may not provide the latest UHD version.

        \item \textbf{Source Code Installation:} Downloads the UHD source code from GitHub and builds it locally. This provides access to the latest features, bug fixes, and allows customization during compilation.
    \end{itemize}
\end{itemize}

\begin{tcolorbox}[
    colback=blue!5,
    colframe=DarkBlue,
    boxrule=0.5pt,
    arc=2pt,
    left=2mm,
    right=2mm,
    top=1mm,
    bottom=1mm
]
\footnotesize
\textbf{Note:} The following steps describe the installation of UHD from \textbf{source code}.
\end{tcolorbox}



\begin{itemize}
    \item \textbf{Step 1: Go to the workspace directory.}
\end{itemize}

\scriptsize
\begin{Verbatim}[xleftmargin=3.8em,breaklines,breaksymbolleft={},breaksymbolright={}]
cd $HOME/git
\end{Verbatim}

\footnotesize

\begin{itemize}
    \item \textbf{\textcolor{DarkBlue}{Step 2: Clone the UHD repository.}}
\end{itemize}

\scriptsize
\begin{Verbatim}[xleftmargin=3.8em,breaklines,breaksymbolleft={},breaksymbolright={}]
# Clone UHD repository
git clone https://github.com/EttusResearch/uhd.git

# Go into UHD directory
cd uhd
\end{Verbatim}

\end{frame}